\documentclass{situationist-preprint}

\title{Fitch's Paradox in Mathematics}
\msid{MTH-12-b}
\msfamily{MTH-12}
\msdiscipline{Philosophy of Mathematics}
\mstype{Consequence}
\mskeywords{Fitch's paradox; knowability; Cartesian restriction; Tennant; incompleteness; anti-realism}

\begin{document}
\maketitle

\begin{abstract}
Fitch's paradox shows that if every truth is knowable, every truth is known. Tennant restricts the knowability principle to Cartesian propositions, those whose being known is consistent. We derive from MTH-12-a that the unknowable arithmetical truth it exhibits is Cartesian, so that, as LOG-19-b holds, the restriction is refuted.
\end{abstract}

\paragraph{The restriction.} LOG-19-a accepts the argument of \citet{Fitch1963}. \citet{Tennant1997} restricts knowability to \emph{Cartesian} propositions, those $p$ for which $Kp$ is consistent. Every Fitch conjunction fails this condition, so, as LOG-19-b grants, the restricted principle survives Fitch's argument.

\paragraph{The consequence.} By Theorem~3.1 of MTH-12-a, some true arithmetical sentence $G$ is not knowable; by Corollary~3.2, $G$ is not of Fitch's form (Definition~2.5). More is true: $KG$ is consistent. $G$ contains no knowledge operator, so $KG$ says only that an arithmetical truth is known, which neither logic nor arithmetic excludes; what excludes it is the effectiveness of our methods (Assumption~2.3 of MTH-12-a). So $G$ is Cartesian, true and unknowable.

\begin{proposition}
Tennant's restricted knowability principle is false.
\end{proposition}

\noindent This is the verdict of LOG-19-b, here derived from the conclusion of MTH-12.

\begin{objection}
An anti-realist will not grant that $G$ is true.
\end{objection}

\begin{reply}
The truth of $G$ follows from the consistency of what is knowable, which the anti-realist also accepts, since knowledge is factive (Assumption~2.4 of MTH-12-a); the reasoning from consistency to the truth of a Gödel sentence is constructive, as Gödel's own proof is \citep{Goedel1931}. The anti-realist must instead deny that our methods are effective, which is the Lucas--Penrose claim that MTH-05 rejects.
\end{reply}

\paragraph{Conclusion.} In mathematics the knowability principle fails even in its restricted form, so Fitch's paradox is not the only obstacle to it.

\begin{thebibliography}{99}\small
\bibitem[Fitch(1963)]{Fitch1963} Fitch, F.~B. (1963). A logical analysis of some value concepts. \emph{The Journal of Symbolic Logic}, 28(2), 135--142.
\bibitem[G{\"o}del(1931)]{Goedel1931} G{\"o}del, K. (1931). {\"U}ber formal unentscheidbare S{\"a}tze der Principia Mathematica und verwandter Systeme I. \emph{Monatshefte f{\"u}r Mathematik und Physik}, 38, 173--198.
\bibitem[Tennant(1997)]{Tennant1997} Tennant, N. (1997). \emph{The Taming of the True}. Clarendon Press.
\end{thebibliography}

\end{document}
